% =====================================================================
\section{Introduction}
\label{sec:intro}
% =====================================================================

More perspectives are assumed to be better: a satellite tile shows layout and
connectivity, a street-level view shows facades, surfaces, and signage, and pairing
them should give a vision--language model a richer picture of a place than either
alone. A growing line of work collects, aligns, and jointly models satellite and
ground views for urban understanding on exactly this premise. But ``should'' is
rarely tested directly---existing benchmarks report one accuracy over a mix of
tasks, some answerable from a single view and some requiring both, so a high number
can mean genuine fusion or simply that one view sufficed. We ask the question the
aggregate hides: \emph{when does the second view actually help?}

We answer it with a controlled audit. We build OELLM, a deterministic,
OpenStreetMap-grounded benchmark that pairs a marked satellite tile with
road-aligned street-level views across \num{40} cities, spanning urban-attribute
tasks (answerable from one location) and cross-view tasks (answerable only from the
relationship between views). Our instrument is a four-way image
ablation---\textsf{blind}, \textsf{sat}, \textsf{sv}, \textsf{full}---from which we
read \emph{fusion synergy} ($\textsf{full}-\max(\textsf{sat},\textsf{sv})$, the gain
of both views over the better single one) and an item-level \emph{interference
asymmetry}. We evaluate at the minimum-sufficient input---one satellite tile plus
one forward street view---so every comparison sits on the same regime, and apply the
audit to a fine-tuned 9B VLM and a panel of raw open VLMs.
\Cref{fig:teaser} previews the setup and the headline effect.

The result is a sharp dissociation. On the seven urban-attribute tasks, the second
view adds only $+\num{0.5}$ points over the better single view on average, and at
the item level it is mildly \emph{harmful}: it overturns a correct answer about nine
times as often as it rescues a wrong one, so a per-item single-view oracle---an upper
bound on single-view routing---beats the two-view model by \num{12} points
($p<\num{e-60}$). On the five cross-view tasks the same model and instrument gain
$+\num{19}$ to $+\num{71}$ points, so the audit registers fusion exactly when the
task requires it. The pattern is robust: it reproduces across raw VLMs of different
families and scales, survives a \num{16}k-token reasoning budget and a multi-turn
self-check, holds when inference-time pruning keeps a single street view (which
retains \num{89}\% of correct answers), and persists when the model is trained with
all four street views instead of one (which adds only $\sim$\num{2} points).

We frame this carefully. It is a \emph{task-dependent} observation about
single-view-decidable attribute recognition at commonly served satellite
resolutions, not a claim that overhead imagery is uninformative or that fusion fails
in general---the cross-view controls show large synergy on the same models. We
report bootstrap confidence intervals on every estimate, audit the benchmark for
text-only leakage, and characterise the items where the second view changes the
answer.

Our contributions are:
\begin{enumerate}[itemsep=1pt,topsep=2pt]
  \item \textbf{A task-dependence result for cross-view fusion.} On seven
  urban-attribute tasks a fine-tuned 9B VLM gains only $+\num{0.5}$ points from a
  second view over the better single view, whereas on five cross-view tasks built
  into the same benchmark it gains $+\num{19}$ to $+\num{71}$---a within-model
  dissociation showing that pairing helps only when the task structurally requires
  it.
  \item \textbf{An item-level interference analysis.} Beyond the small average, the
  second view is net-harmful on the urban tasks: it overturns a correct answer about
  nine times as often as it rescues a wrong one, and a per-item single-view oracle
  exceeds the two-view model by \num{12} points---evidence that the model
  adjudicates poorly between redundant views rather than failing to read them.
  \item \textbf{OELLM, a deterministic paired-view benchmark} of \num{12}
  OSM-grounded task families over \num{40} cities, with a four-way image ablation
  and built-in cross-view controls that make the single-view-sufficient versus
  fusion-required distinction measurable.
  \item \textbf{A reusable audit protocol}---fusion synergy, vision contribution,
  and interference asymmetry---for deciding, before committing to dual-view
  collection, whether a second view is warranted; we show the verdict is robust
  across model families and scales, inference-time view pruning, and training-time
  view count.
\end{enumerate}

% =====================================================================
\section{Related Work}
\label{sec:related}
% =====================================================================

\subsubsection{Remote-sensing vision-language models are single-view.}
A wave of instruction-tuned VLMs target remote sensing, including
GeoChat~\cite{geochat2024}, SkySenseGPT~\cite{skysensegpt2024},
LHRS-Bot and LHRS-Bot-Nova~\cite{lhrsbot2024,lhrsbotnova2025},
EarthGPT~\cite{zhang2024earthgpt}, and the multi-sensor generalist
EarthDial~\cite{earthdial2025}. These models are overhead-only: each consumes a
single aerial or satellite image, and none ingests ground-level imagery
(\cref{tab:models}). Where ``multi-modal'' appears here it refers to multiple
\emph{overhead} sensors---optical, SAR, infrared---rather than to pairing satellite
with street-level views. The dominant suites follow suit: RSVQA~\cite{lobry2020rsvqa},
VRSBench~\cite{li2024vrsbench}, and GeoChat's own benchmark pose questions over a
single overhead image. None isolates the marginal value of a second, ground-level
view. Architecturally these build on LLaVA-style designs~\cite{llava_2023,llava15_2024}
with CLIP or SigLIP encoders~\cite{clip_2021,siglip_2023}; we use four as zero-shot
baselines and discuss their training-data provenance in \cref{sec:setup}.

\subsubsection{Multi-view urban benchmarks.}
A smaller set of benchmarks pairs overhead and ground views and notices the surface
phenomenon we study, but stops short of explaining it. UrBench~\cite{urbench2025}
reports that multi-image and single-image performance differ little on
coordinate-aligned street/satellite pairs, but attributes the gap to question
\emph{comprehension} and measures only aggregate accuracy, leaving open whether the
second view is unused, redundant, or actively harmful. CityLens~\cite{citylens2026},
concurrent with our work, finds street view alone competitive with both views---but
on a three-task socioeconomic \emph{regression} target, with a single model and no
item-level or cross-view analysis. Neither isolates \emph{where in input-mode space}
the effect lives. We contribute the missing instrument: a four-way image ablation
yielding a per-item interference asymmetry and a single-view oracle, applied across a
task taxonomy whose built-in cross-view controls confirm the same instrument
registers large fusion when fusion is required. The result is not ``a second view
rarely helps''---already suspected---but a \emph{mechanism} (the model overturns
about nine correct answers for each it rescues) and a \emph{boundary} (fusion is
decisive exactly when the answer lives between views). Convergent prior observations
make this audit timely; they do not make it redundant. UrbanLLaVA~\cite{urbanllava2025}
trains over satellite, street-view, OSM, and trajectory inputs but poses its
attribute tasks from satellite alone; CityCube~\cite{citycube2026} targets cross-view
spatial reasoning rather than attribute recognition.

\subsubsection{Cross-view geo-localization: where both views are necessary.}
Satellite--ground pairing is well established in cross-view geo-localization, where
the goal is to match a ground photo to an overhead
reference~\cite{durgam2024crossview,zhai2017predicting,zhu2021vigor}. There the two
views are jointly \emph{necessary}, and methods succeed by modeling the
distinct-but-linkable layout they share~\cite{zhang2023geodtr}. This is the regime
our cross-view controls occupy; complementarity that is constitutive of localization
is largely absent from single-view-decidable attribute recognition.

\subsubsection{The marginal value of a second modality.}
That adding a modality need not help is known beyond geosensing. Wang
\etal~\cite{wang2020multimodal} show the best single-modal network can outperform a
naively fused one, and Liang \etal~\cite{liang2023quantifying} formalize multimodal
interaction into redundancy, uniqueness, and synergy, showing low-synergy regimes are
real; our ``fusion synergy'' is the behavioral analogue of their synergy term.
Against this, multimodal learning can provably lower achievable
risk~\cite{huang2021makes}, and several aerial$+$ground systems report recognition
gains~\cite{workman2017unified,machado2021airound}. Most directly counter to our
setting, Suel \etal~\cite{suel2021multimodal} find fusing satellite and street
imagery improves \emph{regression} of income and deprivation; we engage this in
\cref{sec:discussion}, noting it differs in task type (regression vs.\ four-way
recognition), labels, and model. Our finding refines rather than contradicts the
``fusion helps'' literature by naming a regime where it does not.

